% Level 2, Week 1 -- Welcome Back, Python (student workbook)
\chapterweek{1}{Welcome Back, Python}{Every value has a type, and f-strings glue values into friendly messages.}

\section{What you will learn}
\begin{itemize}[leftmargin=1.2em]
  \item The four basic \textbf{types} of value: whole numbers, decimal numbers, text and True/False.
  \item How to build messages with \textbf{f-strings}.
  \item How to turn typed answers into numbers with \texttt{int()}.
  \item How to make smarter choices with \texttt{and}, \texttt{or} and \texttt{not}.
\end{itemize}

\section{Welcome back: what you already know}
From Level 1 Python you know \texttt{print}, variables, \texttt{for i in range(\ldots)}, \texttt{input} with \texttt{int},
\texttt{if}/\texttt{elif}/\texttt{else}, \texttt{random}, turtle drawing, and your own commands made with \texttt{def}.
Level 2 builds on all of them. Two new things will appear in every starter file from today:
\begin{itemize}[leftmargin=1.2em]
  \item Functions that \textbf{hand back an answer} with \texttt{return}, like \texttt{return "Ready to play"}. You will learn all about \texttt{return} next week;
        for now, read it as \enquote{this is the function's answer}.
  \item A line \texttt{if \_\_name\_\_ == "\_\_main\_\_":} near the bottom. The code under it is the main program: it runs when you press \menu{Run}.
        The functions above it are the tools it uses. Keep this line as it is.
\end{itemize}

\section{The big idea: every value has a type}
Python sorts every value into a \textbf{type}. The type decides what Python can do with it.
Here are the four types you will use most:

\begin{center}
%% alt: Four boxes for the four types of value. An int is a whole number such as 7, -2 and 100. A float is a decimal number such as 3.5, 0.0 and 2.0. A str is text in quotes such as "cat" and "7". A bool is True or False.
\begin{tikzpicture}[bin/.style={draw=black,thick,rounded corners=3pt,minimum width=3.5cm,minimum height=2.5cm,align=center,font=\small}]
  \node[bin,fill=kcorange!25] (i) at (0,0) {\textbf{int}\\whole number\\[4pt]\texttt{7}\quad \texttt{-2}\quad \texttt{100}};
  \node[bin,fill=kcsky!25] (f) at (4,0) {\textbf{float}\\decimal number\\[4pt]\texttt{3.5}\quad \texttt{0.0}\quad \texttt{2.0}};
  \node[bin,fill=kcgreen!20] (s) at (8,0) {\textbf{str}\\text in quotes\\[4pt]\texttt{"cat"}\quad \texttt{"7"}};
  \node[bin,fill=kcyellow!35] (b) at (12,0) {\textbf{bool}\\True or False\\[4pt]\texttt{True}\quad \texttt{False}};
\end{tikzpicture}
\end{center}

\begin{tip}[Math helper: numbers below zero]
Numbers can go below zero, like a thermometer on a very cold day. Python writes them with a minus sign in front:
\texttt{-2} means \enquote{2 below zero}. Say it as \enquote{minus two}. You met these in Level 1 turtle drawings:
\texttt{goto(-200, 0)} is 200 steps left of the middle. A number below zero is always smaller than 0.
\end{tip}

\textbf{A float always has a decimal point.} Python shows \texttt{2.0}, not \texttt{2}, when the value is a float.
So \texttt{4 * 12.5} gives \texttt{50.0}: a whole amount, but still written as a decimal number.

\textbf{Quotes are the big clue.} \texttt{7} is a number, but \texttt{"7"} is text that happens to show a 7.
You can ask Python for the type of anything. Type this in the Thonny \textbf{Shell} (after \texttt{>>>}):
\begin{pycode}[Try it in the Shell]
type(7)        # <class 'int'>
type("7")      # <class 'str'>
type(3.5)      # <class 'float'>
type(True)     # <class 'bool'>
\end{pycode}

\subsection{Bug hunt: text plus a number}
Try \texttt{"5" + 5} in the Shell. Python shows a red \textbf{error message}. Always read it from the \textbf{bottom line up}:
\begin{outputcode}
TypeError: can only concatenate str (not "int") to str
\end{outputcode}
\enquote{Concatenate} means \enquote{glue together}. Python can glue text to text, but not text to a number.
This matters because \textbf{\texttt{input()} always gives back text}, even when you type a number. Turn it into a whole number with \texttt{int()}:
\begin{pycode}
level = int(input("Level? "))     # the player types 4
print(level + 1)                  # 5
\end{pycode}

\subsection{f-strings: fill in the blanks}
Put an \texttt{f} in front of the quotes, and Python fills every pair of curly braces with a value.
It works with any type, so you never have to worry about text plus a number.

\begin{center}
%% alt: The f-string f"{name} is level {level}" with name set to "Nova" and level set to 4 fills in to give the text Nova is level 4.
\begin{tikzpicture}[font=\small]
  \node[draw,rounded corners=2pt,fill=kcsky!15,inner sep=5pt] (t) {\texttt{f"\{name\} is level \{level\}"}};
  \node[draw,rounded corners=2pt,fill=kcgreen!15,inner sep=5pt,right=1.6cm of t] (r) {\texttt{Nova is level 4}};
  \draw[-{Stealth[length=3mm]},thick] (t) -- node[above]{fill in} (r);
  \node[below=0.25cm of t,align=center] {\texttt{name = "Nova"}\quad \texttt{level = 4}};
\end{tikzpicture}
\end{center}

\subsection{Two new comparison signs}
In Level 1 you compared with \texttt{==}, \texttt{!=}, \texttt{<} and \texttt{>}. Two more signs save you work.
Each one is typed as two keys next to each other, with no space.
\begin{center}
\renewcommand{\arraystretch}{1.3}
\begin{tabular}{@{}llll@{}}
\toprule
\textbf{Sign} & \textbf{Say it as\ldots} & \textbf{Example} & \textbf{True when\ldots} \\ \midrule
\texttt{>=} & more than or equal to & \texttt{level >= 3} & level is 3 or more (3, 4, 5 \ldots) \\
\texttt{<=} & less than or equal to & \texttt{energy <= 2} & energy is 2 or less (2, 1, 0 \ldots) \\ \bottomrule
\end{tabular}
\end{center}
Every comparison gives back a bool: \texttt{True} or \texttt{False}.

\subsection{Smarter choices: and, or, not}
\begin{center}
\renewcommand{\arraystretch}{1.3}
\begin{tabular}{@{}lll@{}}
\toprule
\textbf{Word} & \textbf{True when\ldots} & \textbf{Game example} \\ \midrule
\texttt{and} & \textbf{both} sides are true & \texttt{has\_key and level >= 3} \\
\texttt{or}  & \textbf{at least one} side is true & \texttt{level >= 10 or has\_map} \\
\texttt{not} & the opposite & \texttt{not game\_over} \\ \bottomrule
\end{tabular}
\end{center}
Each side needs a full comparison: write \texttt{energy >= 5 or energy == 1}, not \texttt{energy >= 5 or 1}.

\begin{tip}[Hero names, not real names]
Your real name, school, address and birthday are \textbf{personal information}. Keep them private.
In our programs you always type a made-up \textbf{hero name} or nickname.
\end{tip}

\begin{exercise}[Hero Card]
Folder: \texttt{exercises/week01/ex-hero-card}. Open \texttt{starter.py} and press \menu{Run} (\keys{F5}).
\begin{enumerate}[leftmargin=1.4em]
  \item It crashes! Read the error from the bottom up. Write the last line here:
  \answerlines{1}
  \item TODO 1 has two parts, and you need both. Part 1: turn the level into a number with \texttt{int()}.
        Part 2: rewrite the \texttt{Next level} line as an f-string, because text \texttt{+} a number still crashes. Run again.
  \item TODO 2--5: add the missing lines with f-strings: power, power points, health (\texttt{level * 12.5}) and \enquote{Legend?}.
  \item Which type is each value? Write int, float, str or bool:\par
  \texttt{hero\_name} \rule{2cm}{0.4pt}\quad \texttt{power\_points} \rule{2cm}{0.4pt}\quad \texttt{health} \rule{2cm}{0.4pt}\quad \texttt{is\_legend} \rule{2cm}{0.4pt}
  \item Done looks like: for level 4 the card shows \texttt{Next level: 5}, \texttt{Power points: 40}, \texttt{Health: 50.0} and \texttt{Legend? False}.
        The whole card is in the folder's \texttt{README.md}.
  \item Check your work: open \texttt{check\_my\_work.py} and press \menu{Run}.
\end{enumerate}
\end{exercise}

\begin{exercise}[Mood Meter]
Folder: \texttt{exercises/week01/ex-mood-meter}. The program asks for your energy (0--10) and whether you had a snack,
then draws a meter. Finish \texttt{mood\_for()}:
\begin{itemize}[leftmargin=1.2em]
  \item \enquote{Super charged}: energy 8 or more \textbf{and} a snack.
  \item \enquote{Ready to play}: energy 5 or more \textbf{or} a snack.
  \item \enquote{Recharge time}: everything else.
\end{itemize}
Before you run it, \textbf{predict} the mood. Then run and check.
\begin{center}
\renewcommand{\arraystretch}{1.5}
\begin{tabular}{|c|c|p{3.5cm}|p{3.5cm}|}
\hline
\textbf{energy} & \textbf{snack?} & \textbf{my guess} & \textbf{the program said} \\ \hline
9 & yes & & \\ \hline
9 & no & & \\ \hline
2 & yes & & \\ \hline
3 & no & & \\ \hline
\end{tabular}
\end{center}
Last step: use \texttt{if not had\_snack:} to print a friendly tip. Then check your work with \texttt{check\_my\_work.py}.
\end{exercise}

\begin{vocab}
  \vocabword{type}{The kind of value: int, float, str or bool. It decides what Python can do with it.}
  \vocabword{int}{A whole number, like 7 or -2.}
  \vocabword{negative number}{A number below zero, written with a minus sign, like -2.}
  \vocabword{float}{A number with a decimal point, like 3.5 or 2.0.}
  \vocabword{str (string)}{Text inside quotes, like \texttt{"cat"} or \texttt{"7"}.}
  \vocabword{bool}{True or False. Comparisons like \texttt{level >= 3} give a bool.}
  \vocabword{f-string}{Text with an f in front. Curly braces inside are filled with values.}
  \vocabword{TypeError}{An error that means \enquote{these two types do not go together}.}
  \vocabword{personal information}{Facts that say who you really are or where to find you. Keep it private.}
\end{vocab}

\begin{reflection}
Which value card or error surprised you today? Why?
\answerlines{2}
Make up a game rule that uses \texttt{and} or \texttt{or}:
\answerlines{2}
\end{reflection}
